\chapter{Inventario normativo}

\begin{cautela}
Este apéndice registra \textbf{lo obtenido, no sólo lo usado}. Catorce de sus entradas ---los Decretos 1.045/96, 1498, 3774/09 y 4.913/98 y las Leyes 3292/58, 4708, 4829/74, 5061, 5114, 5814, 6133, 6895, 7460 y 27.424--- no se citan en ningún capítulo: se consignan porque su texto se consiguió y porque forman parte del marco que quien continúe este trabajo va a necesitar. Que una norma esté aquí no significa que el libro la haya usado; significa que está disponible.

\textbf{Sobre la numeración de las leyes provinciales.} El Digesto Jurídico (Ley 7913, de 2015) renumeró las leyes anteriores. Este libro cita cada ley por el número con que la publica su fuente; cuando conoce los dos, da el del Digesto seguido del original entre paréntesis ---Ley 1423 (orig.\ 145), Ley 2616 (orig.\ 1338)---. Los números de cinco cifras de las leyes de 1929 ---11078 y 11088--- no son un error: hasta comienzos de los años treinta, leyes y decretos se registraban en una sola numeración correlativa, y así los imprime el Boletín.
\end{cautela}

\label{cap:normativa}
\section*{Municipio de La Caldera}

\begin{longtable}{@{}p{2.6cm}p{2.2cm}p{8.2cm}@{}}
\toprule
\textbf{Norma} & \textbf{Fecha} & \textbf{Objeto y estado} \\
\midrule
\endhead
\textbf{Código de Ordenamiento Territorial*} & 2016 (Ord.\ 06/2016) & 9 títulos, 152 artículos, 8 anexos, atlas de 8 planos. \textbf{Texto obtenido} en el sitio del municipio, unas 150 páginas. Declara \textbf{no urbanizables ambas márgenes del río} (zona PU, art.\ 43) \textbf{sin fijar el ancho}, que remite al plano del Anexo II. \textbf{El plano se publica sólo como imagen a escala de 4 km, no mensurable.} Cita su límite jurisdiccional de dos maneras ---leyes 2131/47 y 4987/74 en el artículo 3; 853/47, 4365/70 y 4987/74 en el 36---, que resultan equivalentes: la 2131 es la 853 renumerada y la 4987 modificó la 4365 (cap.~\ref{cap:cot}). \\
\textbf{Ley 5602/80 y Dcto.\ 924/80} & 1980 & \textbf{No obtenidos.} Régimen de \textbf{clubes de campo}. Identificados por Polliotto \textit{et al.} (UCASAL, 2019) como la normativa aplicable a esa figura, distinta de la que rige las urbanizaciones privadas. Bajo este régimen se aprobaron Las Vertientes (2014) y Chacras del Gallinato (2015). \\
\textbf{Res.\ Junta de Catastro 25.819/97} & 1997 & \textbf{No obtenida.} Norma que rige las \textbf{urbanizaciones privadas}, según la misma fuente. Es de rango inferior a la ley y al decreto que regulan los clubes de campo. \\
\textbf{Decreto 14.205*} & 9/12/1931 & \textbf{Texto obtenido.} Interventor Nacional Mariano Gómez, en acuerdo de ministros. Art. 1: fija la categoría de los municipios existentes; \textbf{Caldera, tercera}. Art. 3: las de tercera se gobiernan por comisión de tres a cinco miembros elegidos por el pueblo, con \textbf{presidente nombrado por el Poder Ejecutivo}. Art. 4: los presidentes duran \textbf{un año}; los intendentes, cuatro. Es el acto que fija la categoría que el municipio conserva en 2022. \\
Ord. 040/82 & 1982 & Vigencia mantenida por el art. 4 del proyecto de ordenanza del PDUA. Rige alturas, materiales, cartelería y área de preservación del casco histórico. \pendiente{texto} \\
Ord. 243/97 & 1997 & Ídem. \pendiente{texto} \\
Ord. 643* & 25/11/2015 & Aprueba el PDUA. Informe final y proyecto de ordenanza obtenidos. Tres zonas, once distritos, polígonos de 15 y 25 puntos. No fija porcentajes de cesión. \\
Ord. 733 & 2021--22 & Objeto no determinado. Prorrogada por Res.\ 01/26 en sesión extraordinaria. \\
Ord. 772/23 & 2023 & Moratoria para impuestos y tasas municipales. \\
Ord. 830/25* & prom.\ 2026 & Regula el ingreso al perilago de Campo Alegre; prohíbe todo cobro de acceso al espejo de agua. \\
Ord. 831 y 832 & prom.\ 2026 & Objeto no determinado. \\
Ord. s/n & --- & Régimen de tributos previsto por el art. 30 inc. 2 del COT. \textbf{Verificar existencia.} \\
Ord. s/n & --- & Régimen de multas previsto por el art. 16 inc. b del COT. \textbf{Verificar existencia.} \\
Ord. s/n & --- & Consejo Asesor de Planeamiento, previsto como facultativo por el art. 23. \textbf{Verificar existencia.} \\
Res. 01--04/26* & 21/01/2026 & Prórroga de la Ord.\ 733; clausura preventiva del ingreso a Solanas; clausura de los campings ``Los Teros'' y ``Papi Lalo''; partidas para obras de defensa. \\
\bottomrule
\end{longtable}

\section*{Municipio de Vaqueros}

\begin{longtable}{@{}p{3.2cm}p{1.9cm}p{7.9cm}@{}}
\toprule
\textbf{Norma} & \textbf{Fecha} & \textbf{Objeto y estado} \\
\midrule
\endhead
Ord. 437/14* & prom.\ 29/12/2014 (Dto.\ 78/14); publ.\ 15/01/2015 & Aprueba el Reglamento de Fraccionamiento de Tierras. Fundada en el art. 21 inc. 38 de la Ley 1349. B.O.\ Nº 19.462; factura Nº 95, importe \$95. El cuerpo publicado no consigna número, fecha de sanción ni promulgación, de la que el art.\ 2 hace depender la vigencia. La ficha web no ofrece enlace «Ver anexo». \textbf{El anexo con el reglamento no se publicó}: remite a consulta en oficinas. Orden de publicación 100045820. \textbf{Texto completo obtenido} del Digesto Jurídico Virtual de Vaqueros (copia firmada con el Decreto 78/14 y versión transcripta): 92 artículos. \textbf{Erratas del original y lo que corresponde:} art.\ 26, «artículo 4° Inc.\ a)» $\rightarrow$ \textbf{art.\ 2º inc.\ a} (definición de urbanización o loteo; el art.\ 4 define jerarquías viales); art.\ 68, «art.\ 71° inc.\ f» en la copia firmada y «61° inc.\ f» en la transcripta $\rightarrow$ \textbf{art.\ 61 inc.\ h} (plan de obras y plazo máximo de cinco años; el art.\ 71 trata del control de Catastro y el 61 inc.\ f de los certificados de servicios). \textbf{Vigencia:} el art.\ 2 de la ordenanza la hace regir desde el día siguiente a la promulgación (30/12/2014), mientras el art.\ 90 del reglamento lo aplica a los trámites iniciados después de su «publicación», que en el Boletín nunca ocurrió; una lectura posible es tomar la publicación del acto aprobatorio (15/01/2015); este libro no la resuelve, y en cualquiera de las dos quedan afuera los trámites iniciados antes de fines de 2014. La Ord.\ 535/18 la modifica «en los artículos referidos a urbanizaciones» sin enumerarlos. \\
Código de Planeamiento Urbano Ambiental* & 26/12/2016\newline (texto de 2026) & Ord.\ 507/16, Anexo I (Decreto 165/16); la Ord.\ 567/20 le atribuye el art.\ 101 de retiros, que según la 607/23 es del Código de Edificación; reformado por la Ord.\ 607/23 (17/11/2023; arts.\ 17 a 19 y plano: zonas Mixto 1, IC y polígono RENABAP) y por la 648/26 (26/01/2026, promulgada por el intendente interino Viveros, cuyo Decreto 486/26 la cita como «648/25»: distritos M2 Especial, M2 Costanera y PU Ribera; el Parque Urbano queda sujeto a un acuerdo con la Provincia, titular de las tierras). 6 capítulos, 19 artículos, 3 anexos. Aplica el PIDUA-VAQ. Ámbito: Vaqueros y Lesser. Art. 14 sobre procesos de gentrificación. \textbf{Texto completo obtenido.} \\
Código de Edificación* & 26/12/2016\newline (texto de 2026) & Ord.\ 507/16, Anexo II; deroga el de la Ord.\ 316/2008. Reformado por las Ords.\ 607/23 y 648/26. \textbf{Texto no consultado.} \\
PIDUA-VAQ* & 12/11/2015 & Norma marco. Ordenanza 474/15, sancionada en octubre de 2015 y promulgada por el Decreto 119/15; elaborado en el Programa de Mejora de la Gestión Municipal BID 1855-OC-AR; la 507/16 aprueba los códigos que lo aplican. Texto obtenido. \\
Ord. 495/16* & 07/07/2016 & Acepta la donación de calles y pasajes (2 ha 4.333,21 m²), ochavas (215,64 m²) y espacios verdes (7.102,65 m²) del «Loteo Cerros de Buena Vista», matrículas 1713 y 1905; planos en Inmuebles, expte.\ 18-28358; expte.\ municipal 990/12. Decreto 147/16. \textbf{Errata del original:} el visto y el decreto citan una «Ordenanza nº 529/14» que no existe con ese número ---en 2014 la numeración va de la 419 a la 437 y la 529 es de 2018---; no se identifica a cuál se refiere. \\
Ord. 502/16 y 503/16* & 20/10/2016 & Pasarela en el acceso a Villa Sara y ciclovía en Lesser; los Decretos 156/16 y 157/16 las promulgan subordinando su ejecución a fondos provinciales o nacionales por falta de recursos propios. \\
Ord. 469/15* & 30/10/2015 & Ejidos urbanos de Vaqueros y de Lesser por coordenadas POSGAR 98; funda la jurisdicción sobre Lesser en el art.\ 98 de la Ley 1349, sin citar ley de límites. Decreto 112/15. \\
Ord. 535/18* & 28/09/2018 & Nuevas urbanizaciones: manual de procedimiento en seis pasos, cesiones del 10, 5 y 10\,\%, garantía de obra, licencia de comercialización, audiencia pública en instalaciones municipales. Modifica parcialmente la 437. \\
\bottomrule
\end{longtable}

\section*{Régimen provincial de tierra fiscal}

\begin{longtable}{@{}p{3.2cm}p{1.9cm}p{7.9cm}@{}}
\toprule
\textbf{Norma} & \textbf{Fecha} & \textbf{Objeto y estado} \\
\midrule
\endhead
Ley 2616 (orig. 1338)* & sanc.\ 24/08/1951\newline prom.\ 27/08/1951 & Autoriza a enajenar tierra fiscal. B.O.\ Nº 4.025 del 03/09/1951. \textbf{Vigente}: Digesto Jurídico (Ley 7913), Anexo I, Nº 18 de la nómina, fs.\ 90. Art. 1 inc. a: adjudicación directa para vivienda familiar. Art. 2: orden de preferencia, encabezado por los actuales ocupantes. Art. 5 (sust.\ por Ley 3590/61): precio con coeficiente de disminución del 30\%. Art. 6: máximo dos parcelas por persona o familia. \textbf{Texto completo.} \\
Decreto 4.414/08* & 21/10/2008 & Emergencia habitacional provincial. Crea el Programa «Comunidad Organizada» (Subsecretaría de Tierra y Hábitat) y remite su reglamentación al Ministerio, que la dicta ocho años después por la Res.\ 82/16. Art.\ 4: excluye a los incursos en ocupación ilegal de inmuebles. B.O.\ Nº 17.974. \textbf{Texto completo.} \\
Decreto 450/18 & 2018 & Encomienda al IPV la inscripción, asignación y preadjudicación de inmuebles fiscales a personas de menores ingresos. \pendiente{texto} \\
Decreto 1.045/96 & 1996 & Invocado en el visto de las Resoluciones 7, 8 y 9 de 2016. \pendiente{texto} \\
Ley 7.905 y Decretos 18/15 y 21/15 & 2015 & Crean el Ministerio, la Secretaría de Tierra y Bienes y la Subsecretaría de Tierra y Hábitat. \\
Res. 82/16* & publ.\ 09/09/2016 & Reglamento de Asignación de Lotes de Interés Social. Firmante: Saravia. B.O.\ Nº 19.861. \textbf{El anexo ---once artículos--- se publica sólo como archivo vinculado, escaneado y sin capa de texto}: copia certificada de seis folios (cap.\ \ref{cap:tierrafiscal}). \\
Res. 86/16* & publ.\ 09/09/2016 & Publicada en el mismo Boletín que la 82/16. Modifica el art. 7º del Anexo de la 82/16: sustituye ``posesión'' por ``goce'' en la definición de tenencia precaria. \\
Res. 7/16* & 17/11/2016 & Requisitos y documentación; formulario FAS001. Incompatibilidad de autoridades superiores (art.\ 2 inc.\ d); exclusión de incursos en ocupación ilegal (inc.\ c); ratificación anual del legajo y baja a los tres años (art.\ 8); inscripción en la Subsecretaría o en el municipio donde se lleve adelante la urbanización, si hay convenio vigente (art.\ 8 inc.\ a); acceso únicamente por sorteo público, salvo asignación directa (inc.\ j); impugnación por la Ley 5.348 (art.\ 9). B.O.\ Nº 19.906; fe de erratas en el Nº 19.911, que alcanza además a una Resolución 3/16 de la misma Secretaría que este libro no tiene. \textbf{Texto del anexo obtenido}: se publica como imagen escaneada vinculada al acto, trece folios ---seis del Anexo I y siete del formulario FAS001 del Anexo II---, foliados 12 a 24 y certificados como copia fiel. \\
Res. 8/16* & 17/11/2016 & Tipologías de puntaje. B.O.\ Nº 19.906; fe de erratas en el Nº 19.911. \textbf{Anexo publicado sólo como imagen escaneada} (tres folios). \\
Res. 9/16* & 17/11/2016 & Procedimientos de asignación; cupos del 70/28/2\%; garantías del sorteo; acta de tenencia precaria (Anexo II, modelo). B.O.\ Nº 19.906; fe de erratas en el Nº 19.911. \textbf{Texto del anexo obtenido}: se publica como imagen escaneada vinculada al acto, trece folios ---diez del Anexo I y tres del modelo de Acta de Tenencia Precaria---, foliados 33 a 45 y certificados como copia fiel. \\
Decreto 399/20* & 24/06/2020 & Crea el Plan ``Mi Lote''. Meta de 10.000 lotes, 40\% metropolitano. Implementación a cargo del IPV. \\
Decreto 1639/07* & 2007 & Aprueba las actas con Vaqueros y La Caldera por el Servicio de Información Catastral de Inmuebles; precio, 30\,\% de la recaudación adicional (15\,\% transitorio en La Caldera); los anexos dejan constancia de que ninguno de los dos municipios tenía código de planeamiento. B.O.\ Nº 17.645 (cap.\ \ref{cap:hacienda}). \\
Decreto 3318/08 & 2008 & Designa la Comisión Directiva de la Cooperadora Asistencial de La Caldera, presidida por el intendente Calabró. B.O.\ Nº 17.931. \\
Decreto 2813/09 & 2009 & Declara de interés provincial la cooperación con Champagne-Ardenne para conformar el Corredor Intermunicipal; modelo de los Parques Naturales Regionales franceses. B.O.\ Nº 18.145. \\
Decreto 4354/09* & 2009 & Aprueba el convenio del Corredor Intermunicipal para el Desarrollo Sustentable (CorInDeS) entre La Caldera, Vaqueros, San Lorenzo, Cerrillos y Salta, en el marco de la cooperación con Champagne-Ardenne; incluye el ordenamiento territorial entre sus fines. B.O.\ Nº 18.207 (cap.\ \ref{cap:redes}). \\
Decreto 2508/09 & 2009 & Aprueba el convenio de factibilidad con el PAMI para refaccionar y ampliar el Hogar de Ancianos de La Caldera; Comité Coordinador con el intendente. B.O.\ Nº 18.129. \\
Decretos 89/15 y 1634/17* & 2015 y 2017 & Designan a Héctor Miguel Calabró secretario de Relaciones Institucionales (B.O.\ Nº 19.684, 17/12/2015) y Secretario de Asuntos Municipales (B.O.\ Nº 20.155, 30/11/2017); el Decreto 1608/17 acepta su renuncia al primer cargo y la Decisión Administrativa 592/18 le asigna el adicional por responsabilidad. \\
Decreto 16/19* & 10/12/2019 & Aprueba las estructuras ministeriales del nuevo gobierno; la Secretaría de Tierras y Bienes del Estado, con la Subsecretaría de Regularización Dominial, queda en la Secretaría General de la Gobernación. B.O.\ Nº 20.649. Anexo en imagen. \\
Decreto 127/20* & 22/01/2020 & Estructura de la Secretaría General, vigente desde el 11/12/2019: la Dirección General de Inmuebles depende de la Secretaría de Tierra y Bienes del Estado. B.O.\ Nº 20.672. Anexo escaneado. \\
Decreto 844/21* & 27/09/2021 & Traslada la Secretaría al Ministerio de Infraestructura; la DGI depende directamente del Ministro. \\
DA 569/25* & 23/12/2025 & Estructura de la Jefatura de Gabinete, vigente desde el 15/12/2025: seis secretarías (entre ellas Obras Públicas y Tierra y Bienes del Estado), Inmuebles bajo el jefe de Gabinete y la Subsecretaría de Recursos Hídricos dentro de Obras Públicas. B.O.\ Nº 22.100. Anexos escaneados. \\
Decreto 63/25* & 06/02/2025 & Adjudica en venta ciento tres parcelas del barrio Juan Manuel de Rosas. \textbf{Anexo nominal obtenido} (B.O.\ Nº 21.894; anexo SA100050047): copia certificada escaneada de cuatro folios, con capa de texto de reconocimiento óptico imperfecta. Verificados sobre el anexo: 103 parcelas en las manzanas 430 a 475, matrículas 188.191 a 188.612, y doce parcelas con trece expedientes de la serie anterior (1995 a 2018). \\
Res. 567D/26 & 2026 & Subsana errores materiales en anexos de dieciocho instrumentos, entre ellos el Decreto 63/25. Anexos I (Capital) y II (Interior) \textbf{publicados como imagen escaneada vinculada} (SA100053553): veinte rectificaciones; dos cambian parcela y matrícula y una convierte el precio en «sin cargo». El anexo titula «159D/23» lo que el texto llama 156D/23 (cap.\ \ref{cap:tierrafiscal}). \\
\bottomrule
\end{longtable}

\section*{Antecedentes provinciales sobre el suelo del departamento}

\begin{longtable}{@{}p{3.2cm}p{1.9cm}p{7.9cm}@{}}
\toprule
\textbf{Norma} & \textbf{Fecha} & \textbf{Objeto} \\
\midrule
\endhead
\textbf{Decreto 2478/12 --- PDES 2030*} & 2012 & \textbf{Texto e Informe Final obtenidos y leídos completos.} Aprueba el Plan de Desarrollo Estratégico de la Provincia de Salta 2030 con «carácter definitorio» según el art.\ 77 de la Constitución Provincial. Elaborado por iniciativa privada de la Fundación Salta, con metodología del IAE (Universidad Austral) y financiamiento del BID/DINAPREI. \textbf{Trece sectores, siete ejes provinciales, sin estructura territorial por departamento.} Declara que el crecimiento del área metropolitana fluyó «espontáneamente, sin planificación», mientras su capítulo de Minería cartografía los salares de la puna con sus empresas \textit{(cap.\ \ref{cap:trabajo})}. \textbf{Documento leído completo}: «La Caldera» aparece dos veces en sus ciento setenta páginas ---en la nómina de participantes del Diagnóstico y en un inventario hotelero---. \\
\textbf{Decreto 1.889/16} & 2016 & \textbf{Decreto y anexo obtenidos.} El decreto se publica como texto en la ficha web (B.O.\ Nº 19.913, 29/11/2016; refrendan Gomeza, Montero y Simón Padrós); el anexo arranca en el folio 3. Contiene el Acta Acuerdo del «Plan de Mínimas y Encauzamiento 2016--2017» entre la Provincia y veintiún municipios, entre ellos La Caldera y Vaqueros. Financiado con la autorización de crédito de la Ley 7.931. Anticipo del 30\,\%, pago contra certificados, fondos a cuentas municipales. \textbf{Los montos por municipio y los proyectos de obra que el acta invoca no constan.} \\
\textbf{Ley 1336*} & sanc.\ 18/08/1951\newline prom.\ 24/08/1951 & \textbf{Ley Orgánica sobre Régimen de Expropiación de la provincia de Salta}, renumerada \textbf{Ley 2614} por el Digesto. B.O.\ Nº 4.025 del 03/09/1951. \textbf{Vigente}: Digesto Jurídico (Ley 7913), Anexo I, Nº 17 de la nómina; el texto empieza a fs.\ 85. Modificada en su art. 34 por la Ley 2661 (original 1383), B.O.\ 11/10/1951. \pendiente{lectura del texto completo en el Digesto} \\
Ley 4272* & 29/11/1968 & \textbf{Texto obtenido.} B.O.\ Nº 8.205 del 11/12/1968. Modifica los arts.\ 13, 15 a 22, 35, 37 y 38 de la Ley 1336: Tribunal de Tasaciones y posesión inmediata consignando la valuación fiscal más 30\% (art.\ 17). Invocada por la Ley 6354/86 (cap.\ \ref{cap:expropiacion}). \\
«Ley 4.172/68» & --- & \textbf{Número erróneo.} Así la cita la Ley 4486 de 1972 para los arts.\ 13 y 17 del régimen de expropiación; la norma de 1968 que reformó esos artículos es la Ley 4272 (cap.\ \ref{cap:expropiacion}). \\
\textbf{Ley 4486*} & 08/06/1972 & Declara de utilidad pública y sujetas a expropiación cinco fracciones del departamento La Caldera para el Embalse de Campo Alegre y el desvío de la RN 9: 686,19 ha en total. Sancionada y promulgada por el Poder Ejecutivo de facto en un solo acto. B.O.\ Nº 9.055 del 19/06/1972. \\
Ley 4708* & 23/11/1973 & Incluye en el Plan de Obras Públicas 1974 el cordón cuneta y pavimentación de la calle principal de La Caldera y la ruta al Cristo Redentor. \\
Ley 6181* & 20/10/1983 & Expropia 1 ha del catastro 132 (Hilario Arias) para la Escuela Nº 660, Finca Mojotoro. Valor fiscal más 30\%: 12,49 pesos argentinos. Dictada bajo facultades de la Junta Militar. \\
Ley 6354* & 04/03/1986 & Expropia 8.158,71 m² de Elena Serrey\index{Serrey, Elena} de González Bonorino (matrícula 1.303, localidad de La Caldera) para viviendas familiares; donación al IPDUV con reversión automática si no se construye en dos años. Tramita por Leyes 1.336 y 4.272. \\
Ley 7781* & 2013 & DNU 1595/13: Emergencia Hídrica en todo el territorio provincial por escasez de agua, atribuida al cambio climático. Ordena readecuar y ordenar todas las cuencas. \\
Ley 4365* & 05/08/1970 & \textbf{Texto obtenido.} B.O.\ Nº 8.612 del 18/08/1970. Crea el municipio de tercera categoría (2da.\ clase) de Vaqueros; límites: al norte, río Wierna; al sur, río Vaqueros (Capital); al este, río La Caldera; al oeste, fincas Lesser y Abra de Lesser, del municipio de La Caldera. Modificada por la Ley 4987 (cap.\ \ref{cap:cot}). \\
Ley 2131 (orig.\ 853)* & 24/07/1947 & \textbf{Texto obtenido.} Promulgada el 02/08/1947 (Lucio A.\ Cornejo); B.O.\ Nº 2.912 del 07/08/1947. Delimita los distritos municipales de la provincia: «El Municipio de la Caldera abarca todo el territorio del Departamento». Art.\ 2: hasta un censo general, los municipios mantienen sus categorías y organización (caps.\ \ref{cap:cot} y \ref{cap:poblacion}). \\
Ley 4987* & 19/09/1974 & \textbf{Texto obtenido.} Promulgada el 17/10/1974 (Ragone); B.O.\ Nº 9.613 del 30/10/1974. Límite entre La Caldera y Vaqueros: río Potrero de Castilla, río Los Yacones y río Wierna hasta su confluencia con el río La Caldera. Modifica las leyes 853/47 y 4365/70. Art.\ 3: deslinde y amojonamiento por la Dirección General de Inmuebles en noventa días (cap.\ \ref{cap:cot}). \\
\bottomrule
\end{longtable}

\section*{Ordenamiento territorial de bosques nativos}

\begin{longtable}{@{}p{3.2cm}p{1.9cm}p{7.9cm}@{}}
\toprule
\textbf{Norma} & \textbf{Fecha} & \textbf{Objeto} \\
\midrule
\endhead
Ley nacional 26.331 & 2007 & Presupuestos mínimos de protección ambiental de los bosques nativos. \\
Ley provincial 7543 & 2008 & Ordenamiento Territorial de Bosques Nativos de Salta. Su art. 14 fue modificado por la Ley 8483. \\
\textbf{Ley provincial 8483*} & B.O.\ 03/01/2025 & Revisión del OTBN. Modifica superficies por categoría; introduce el Área de Producción y Conservación (APC) y el ``área verde'' sin ubicación geográfica inicial; habilita ganadería silvo-pastoril en Categoría II; prevé certificación de bonos verdes. Anexo I: informe técnico del equipo convocado por Decreto 3749/2014 (SAyDS, INTA, INENCO-CONICET, UNSa, INAI, APN), con taller participativo vinculante del 25/10/2023, Consejo Asesor de 53 instituciones y 91 participantes; incluye el Anexo IX de mapas --- entre ellos el Mapa 11 de comunidades indígenas y campesinas --- y el Anexo X de superficies y mapas por cuenca. El Anexo I cierra con una \textbf{adenda} de la autoridad provincial que responde a las observaciones de la autoridad nacional, ajusta las superficies y repite el porcentaje categorizable por cuenca (Mojotoro--Lavayén: 29.744 ha de APC, 0,0\,\%; sombreada como excedida en la Tabla 17). Anexo II: cartografía aprobada. \textbf{Es la norma votada en contra por el senador por La Caldera en diciembre de 2024.} Su articulado: \textbf{art.\ 5º} ---sustituye el art.\ 20 de la Ley 7543--- exige la \textbf{no objeción de la autoridad forestal} para urbanizaciones y obras con cambio de uso de suelo en zonas del OTBN; \textbf{art.\ 3º}, consulta libre, previa e informada conforme Convenio OIT 169; \textbf{art.\ 9º}, plazo de 180 días ---vencido--- para identificar los desmontes ejecutados y su legalidad; \textbf{art.\ 7º}, sin efecto retroactivo. \\
Res. SAyDS 00509* & B.O.\ 27/08/2025 & Inmuebles con APC pueden solicitar cambio de uso identificando en el APC la Categoría III; el remanente queda en Categoría II. \\
Disp.\ nacional 2483/2025* & B.O.\ 20/10/2025 & Acredita la actualización. Cat.\ I 1.308.244 ha; Cat.\ II 5.528.753 ha; Cat.\ III 721.568 ha; total 7.558.565 ha. \\
Res. SAyDS 00235* & B.O.\ 13/04/2026 & Reglamento de requisitos para proyectos de intervención sobre bosques nativos. Deroga las Res.\ 558/13, 831/19, 333/20 y 411/20. \\
\textbf{Ley provincial 7107*} & sanc.\ 12/10/2000; B.O.\ 08/11/2000 & \textbf{Texto completo obtenido.} Sistema Provincial de Áreas Protegidas, dictada en cumplimiento del art. 98 de la Ley 7070. \textbf{No contiene artículo derogatorio}; su art. 63 remite supletoriamente a la Ley 7070 y su art. 35 inc. c faculta a recategorizar las áreas existentes. Art. 5 inc. e: conservar \textbf{a perpetuidad} los ambientes que circundan las nacientes de cursos de agua; inc. j: proteger zonas de alto riesgo de desastres naturales. Art. 17: once categorías. Art. 26: reservas municipales y convenios de \textbf{transferencia de terrenos fiscales} a municipios. Art. 47 inc. 2: beneficios impositivos y reducción de tasas municipales para propietarios que adhieran. Art. 51: la Autoridad \textbf{debe impulsar expropiación} si las restricciones impiden el ejercicio regular de los derechos. Art. 52: declaración por \textbf{acta multilateral} Provincia--Municipios, inalterable para las signatarias. Art. 53: adhesión voluntaria por tiempo indeterminado, renuncia no antes de veinte años. Art. 57: las restricciones \textbf{se anotan en los Registros} y siguen al inmueble; la D.G.I.\ debe comunicar toda transferencia. Arts. 60--61: la declaración se hace por \textbf{acto administrativo del Poder Ejecutivo}, de oficio o a pedido de parte --- \textbf{no requiere ley}. Art. 62: declarada por decreto, sólo puede alterarse \textbf{por ley}, bajo pena de nulidad.\\
\bottomrule
\end{longtable}

\section*{Régimen provincial de aprobación de loteos}

\begin{longtable}{@{}p{3.2cm}p{1.9cm}p{7.9cm}@{}}
\toprule
\textbf{Norma} & \textbf{Fecha} & \textbf{Objeto} \\
\midrule
\endhead
Ley 2.308 (orig.\ 1.030) & --- & Ley de Catastro. Su Título V rige la creación de nuevos centros de población y la ampliación o modificación del trazado de los existentes. \\
\textbf{Ley 8126*} & sanc.\ 13/11/2018; B.O.\ 17/12/2018 & Régimen de Municipalidades para los municipios sin carta orgánica. Sin categorías ni comisiones municipales; ausencia e impedimento del intendente a cargo del presidente del Concejo (arts.\ 63--64); acefalía definitiva con más de un año de mandato: elecciones (art.\ 67); registro de ordenanzas de consulta pública y publicación en el Boletín Oficial y en internet (arts.\ 78--79). Deroga la Ley 1349 (art.\ 93). \textbf{Texto completo consultado.} \\
\textbf{Ley 1349*} & 1933; derog.\ 2018 & \textbf{Ley Orgánica de Municipalidades, derogada por el art.\ 93 de la Ley 8126. Texto completo obtenido} (versión del Compendio Electoral 2019). Art. 2: la categoría se fija por ley. Art. 3: comprobado un censo, el P.E.\ envía el proyecto de recategorización. Art. 4: primera categoría $>$10.000 hab.; segunda $>$5.000. Art. 5: Concejo de 9 y de 5 concejales respectivamente. Art. 21 inc. 38: los Concejos dictan planes reguladores y disponen la extensión del ejido, conforme el Capítulo XIII. Art. 25: en caso de renuncia del intendente lo suceden los presidentes o vicepresidentes del Concejo. Art. 32--33: tercera categoría = Comisión Municipal, presidente designado por el P.E. Art. 68: los conflictos los dirime la Corte de Justicia. Arts. 69--74: intervención. Art. 79: tope del 25\% de la renta en sueldos administrativos. Art. 80: 10\% de la renta a educación común. \textit{El capítulo \ref{cap:hacienda} registra que en 1933 esos topes se invocaban como artículos de la Ley 68. La identidad está confirmada por el articulado: el Decreto 1498 G del 10 de marzo de 1941 (B.O.\ Nº 1.897) aprueba el ejido de Metán «de acuerdo a lo que establece el art.\ 98 de la ley N.o 68 --- Orgánica de Municipalidades», y el artículo 98 de la Ley 1349 es la norma que manda fijar el ejido por ordenanza con aprobación del Poder Ejecutivo. La tabla oficial de leyes con doble numeración lo confirma: la Ley 1349 es la «original 68», publicada el 3 de marzo de 1933.} \textbf{Cap.\ XIII (arts.\ 97--103): régimen municipal de loteo.} Art. 98: el ejido se fija por ordenanza con aprobación del P.E. Art. 102: el loteador entrega un plano a la Municipalidad y copia a Rentas; sólo requiere autorización previa si proyecta calles nuevas o pasajes. Art. 103: para pueblos nuevos, aprobación previa del plano y fijación del ejido. \\
\textbf{Decreto 1.410/1973*} & B.O.\ 9.371, 24/10/1973 & \textbf{Texto completo obtenido.} Reglamentación para la aprobación de planos de loteos urbanos y suburbanos. Vigente 46 años; derogado por el Decreto 1682/19. Deroga el Decreto 8.369/70. Art. 1: anteproyecto \textbf{visado por la Municipalidad}. Art. 3: exige \textbf{curvas de nivel con equidistancia de 1 m}, zona de reserva para uso público y anteproyecto de obras y servicios. Art. 4: memoria con los servicios que el propietario se compromete a incorporar. Arts. 5--8: resuelven Arquitectura, Viviendas, \textbf{Administración General de Aguas} y Vialidad. Art. 11: aprobación \textbf{por decreto del P.E.} Art. 12: \textbf{previo a la venta el propietario debe cumplir las obligaciones contraídas}. Art. 13: excepciones --- planes oficiales de vivienda o \textbf{fianza}. Art. 14: \textbf{el cumplimiento lo verifica la Municipalidad}. Art. 15: recién entonces la D.G.I.\ autoriza la venta y \textbf{asigna nomenclatura catastral y matrícula}. \\
Res.\ Jefatura de Gabinete 21/2018 & 2018 & Crea una Comisión para tratar la problemática de loteos, urbanizaciones, clubes de campo y asentamientos. \\
\textbf{Decreto 1682/19*} & B.O.\ 28/11/2019 & Aprueba la reglamentación del Título V de la Ley 2.308. Crea, por su art. 11, la Unidad Coordinadora de Loteos, Loteadores y Desarrolladores Inmobiliarios. \textbf{Anexo publicado como archivo vinculado, con texto legible y firma digital}: 33 artículos en seis capítulos (procedimiento, Unidad Coordinadora, preventa, matrícula anticipada, penalidades, transitorias). Exige certificados de no inundabilidad condicionado y definitivo, certificado de aptitud ambiental y caución por las obras; admite la preventa autorizada y la matrícula anticipada (cap.\ \ref{cap:loteo}). \\
Res.\ 466 D/20* & 2020 & Encomienda a la \textbf{Secretaría de Tierras y Bienes} las competencias de la Unidad Coordinadora de Loteos. \\
Decreto 308/22* & 2022 & Delega en el Ministro de Infraestructura la facultad de aceptar donaciones de inmuebles para uso institucional en el marco de la Ley 2.308. \\
Res.\ Conjunta 001/24 ENRESP, 173/24 SRH y 683/24 SAyDS* & 24/09/2024 & Habilita un período extraordinario de regularización del servicio de tratamiento de efluentes cloacales operado de hecho por particulares en área no servida por COSAYSA. \\
Res.\ ENRESP 95/2026* & 14/01/2026 & Procedimiento para la tramitación de factibilidades de agua potable, colección cloacal y tratamiento de efluentes. \\
\bottomrule
\end{longtable}

\section*{Régimen nacional de expansión de redes de gas}

\begin{longtable}{@{}p{3.2cm}p{1.9cm}p{7.9cm}@{}}
\toprule
\textbf{Norma} & \textbf{Fecha} & \textbf{Objeto} \\
\midrule
\endhead
Decreto 2452/92* & 1992 & Otorga a la actual Naturgy NOA S.A.\ (ex Gasnor S.A.) la licencia de distribución de gas natural sobre Salta, Jujuy, Tucumán y Santiago del Estero, por 35 años prorrogables. \\
Res. ENARGAS I-910/09* & 09/10/2009 & Procedimiento para la expansión de sistemas de distribución. Admite obras a instancia de la distribuidora, subdistribuidor o terceros interesados. Define el Valor de Negocio: si supera la inversión, la obra corresponde a la distribuidora; si no, ésta aporta ese monto y la diferencia la cubren usuarios o terceros. Deja sin efecto las Res.\ 10/93 y 44/94. \\
Res. ENARGAS 630/2019* & 2019 & Modalidades de contraprestación a terceros interesados; para usuarios residenciales, exclusivamente en metros cúbicos de gas. \\
Res. ENARGAS 778/2025* & B.O.\ 21/10/2025 & \textbf{No modifica el régimen}: pone en consulta pública, por quince días hábiles, el proyecto de reglamentación del art. 16 de la Ley 24.076 --- T.O.\ 2025 (Anexo IF-2025-115674567). Expte.\ EX-2025-112087756-APN-GD\#ENARGAS. \\
Res. ENARGAS 875/2025* & 06/11/2025 & Amplía el plazo de la consulta por diez días hábiles más. \\
\textbf{Res. ENARGAS 435/2026*} & B.O.\ 27/04/2026 & \textbf{Norma vigente.} Aprueba la nueva Reglamentación de las autorizaciones del art. 16 de la Ley 24.076 --- T.O.\ 2025 (Anexo IF-2026-41276028). Su art. 2º \textbf{deja sin efecto la Res.\ I-910/09} y su normativa concordante. Reduce el horizonte de valuación de 35 a 10 años; reemplaza costos medios por costos marginales de declaración jurada, actualizados por 50\% IPC + 50\% IPIM; impone aplicativo web obligatorio desde la factibilidad. Los trámites iniciados antes siguen bajo el régimen anterior. \textbf{Entre los observantes de la consulta figura Naturgy NOA S.A.} (27/11/2025). \\
Ley 27.424 & 2017 & Régimen de fomento a la generación distribuida de energía renovable integrada a la red. \\
\bottomrule
\end{longtable}

\section*{Designaciones}

Decreto 120/15 (Martinis Mercado, Secretario de Tierra y Bienes); Decisión Administrativa 50/21 (Puertas, Jefe de Programa de Planificación Urbana); Decisión Administrativa 366/22 (Haro, Director General de Tierras Fiscales); Decreto 64/23 (Abilés, Secretaria de Tierras y Bienes del Estado); Decisión Administrativa 19/23 del 26/12/2023 (B.O.\ Nº 21.621: estructura del Ministerio de Infraestructura, con Tierra y Bienes e Inmuebles; redesigna a Haro y Puertas; no confundir con la DA 19/23 de enero de 2023, contrato de servicios de Desarrollo Social); Decretos 197/24 y 198/24 (intervención de PRO.VI.PO.\ S.E.\ y autorización para suscribir actos jurídicos).



\section*{Contrataciones y régimen de aguas}

\begin{longtable}{@{}p{3.2cm}p{1.9cm}p{7.9cm}@{}}
\toprule
\textbf{Norma} & \textbf{Fecha} & \textbf{Objeto} \\
\midrule
\endhead
\textbf{Decreto 1.989/02} & 28/10/2002 & \textbf{Texto completo obtenido.} Aprueba la Resolución 070/02 de la Agencia de Recursos Hídricos, reglamentaria del art.\ 126 del Código de Aguas. Publicado como \textbf{separata} del B.O.\ del 11/11/2002. Anexo I: procedimiento en doce artículos. Anexo II: conceptos y denominaciones. Firmas: Romero, Gambetta, David; la resolución, por el Ing.\ Gustavo López Asensio. \\
Ley 8.072 y Decreto 1319/18* & --- & Sistema de Contrataciones de la Provincia. Su art. 15 enumera los supuestos de \textbf{contratación abreviada}. El inc. i) es el de razones de urgencia; el \textbf{inc. e)} --- reglamentado por el art. 20 del Decreto 1319/18 --- exige fundar la capacidad científica, técnica, tecnológica, profesional o artística de la persona o empresa a contratar. Es el encuadre invocado para la consultoría del Plan de Manejo Integral (cap. \ref{cap:amparo}). \pendiente{texto del art. 15}\\
\textbf{Código de Aguas (Ley 7017)* y Decreto 2299/03*} & 21/12/1998 & \textbf{Texto íntegro obtenido}, con el veto parcial del Decreto 4.913/98 (arts.\ 6, 51, 56 y 310), que la promulgó el 24/12/1998; publicada en el B.O.\ Nº 15.569 del 11/01/1999. En el texto vigente que publica el Digesto de Diputados, «conforme veto parcial», \textbf{el art.\ 6 conserva la creación del Consejo Asesor Provincial del Agua} integrado por los usuarios, y el art.\ 51 conserva la publicación de las modificaciones del Catastro de Aguas en el Boletín Oficial y en un diario de circulación provincial. El Consejo funciona: según la Secretaría de Prensa (07/12/2010), sesiona el primer lunes de cada mes en la Secretaría de Recursos Hídricos, lo integran representantes de seis zonas de consorcios ---La Caldera, en la del Valle de Lerma--- y la Asociación de Consorcios de Usuarios de Agua Pública. El Decreto 3774/09 (B.O.\ Nº 18.182 del 04/09/2009; el 4746/09 sumó Molinos a los Valles Calchaquíes) fijó esas zonas, lo puso bajo la presidencia del secretario de Recursos Hídricos, con doble voto, y exigió quórum de dos tercios; sus dictámenes no son vinculantes. Se obtuvo también la reglamentación completa. Arts.\ 24, 45/47, 51, 62/68, 140/158, 201/215 y 309: concesiones, registro y procedimiento. \textbf{Art.\ 143: las concesiones de aguas subterráneas son eventuales.} \textbf{Art.\ 66: los clubes de campo se consideran ``población''.} \textbf{Arts.\ 46, 258 y 259: obligación de aforar antes de conceder dotación permanente} \textit{(cap.\ \ref{cap:aguabaja})}. \textbf{Art.\ 170: permiso previo para talar en las márgenes.} \textbf{Arts.\ 171 y 172: evaluación previa a concesiones de áridos, y zonas inundables} --- la reglamentación del 172 fija el Certificado de No Inundabilidad como \textbf{recaudo indispensable} y arancelado \textit{(cap.\ \ref{cap:ribera})}. \textbf{Art.\ 236 de la reglamentación: autorización previa para obras de encauzamiento} \textit{(cap.\ \ref{cap:amparo})}\\
\textbf{Ley 7017, art. 126*} & --- & \textbf{Texto obtenido.} Determinación de la línea de ribera conforme al criterio del \textbf{art. 2.577 del Código Civil} (derogado en 2015; el CCyC, art.\ 235 inc.\ c, remite al promedio de las máximas crecidas ordinarias) y al procedimiento que fije la reglamentación, \textbf{dando intervención en la operación a los interesados}. Anotación en el Libro de Catastro. \textbf{In fine: la Autoridad puede rectificar la línea cuando el cambio de circunstancias lo haga necesario.} \\
\textbf{Ley 7017, art. 172*} & --- & \textbf{Texto obtenido} (transcripto en el edicto de la Res. 348/13). \textbf{Zonas inundables.} La Autoridad determina, en las zonas ribereñas, los sectores afectables por inundaciones, \textbf{en los cuales no se permitirá la instalación de asentamientos poblacionales} ni la fijación de obstáculos al curso de las aguas. \textbf{Las nuevas construcciones o plantaciones en zonas ribereñas requieren autorización previa de la Autoridad}, debiendo proyectarse y ejecutarse las obras necesarias para prevenir el riesgo de inundación. \\
\textbf{Ley 7017, arts. 171 a 173*} & --- & \textbf{Textos obtenidos.} Art. 171, información previa: antes de otorgar permisos de uso de cauces, márgenes o \textbf{extracción de áridos}, la Autoridad evalúa el efecto sobre riberas y flujo; si es desfavorable, desaconseja el otorgamiento o exige obras y técnicas de operación. Art. 172, zonas inundables (ver cap. \ref{cap:ribera}): \textbf{sin plazo}, a diferencia del art. 194 del código cordobés, que fija diez años. Art. 173, penalidades: multa, sanciones conminatorias y \textbf{potestad de ordenar la demolición de las obras} o destruirlas por cuenta del infractor. \\
\textbf{Reglamentación, art. 172º*} & --- & \textbf{Texto obtenido.} La Autoridad de Aplicación expide las certificaciones de ``no inundabilidad'', que \textbf{constituyen un recaudo indispensable para todo asentamiento poblacional, urbanización o plantación en zona ribereña}, previo cobro de arancel. \\
\textbf{Ley 7017, art. 143*} & --- & \textbf{Texto obtenido.} \textbf{Todas} las concesiones de uso de aguas subterráneas son \textbf{eventuales}, salvo las obtenidas por declaración judicial de aguas privadas. \\
Res. 070/02 & 2002 & Reglamentación del art. 126. Según el edicto del río Bermejo (B.O. 20.457, 07--08/03/2019): resolución previa que aprueba la \textbf{Comisión Técnica} con sus integrantes nominados, declaración del \textbf{método} a emplear, delimitación de la zona, \textbf{prohibición cautelar} a los ribereños de modificar las márgenes, publicación por dos días y \textbf{quince días hábiles administrativos} para que terceros aporten estudios. Texto obtenido como parte del Decreto 1.989/02.\\
\bottomrule
\end{longtable}

\section*{Otros instrumentos citados en el cuerpo}

{\small\itshape Se consignan con la fecha y el capítulo donde se analizan. Su texto se obtuvo en la medida en que ese capítulo lo transcribe.\par}

\begin{longtable}{@{}p{3.2cm}p{1.9cm}p{7.9cm}@{}}
\toprule
\textbf{Norma} & \textbf{Fecha} & \textbf{Objeto y capítulo} \\
\midrule
\endhead
Ley 1423 (orig.\ 145) & sanc.\ 02/10/1934\newline prom.\ 06/10/1934 & Región económica del olivo; exención de diez años. Publicada el 12/10/1934. Figura como vigente en el Anexo I del Digesto (Ley 7913), cuya nómina la rotula, por error, «Convenio límites interprovinciales»; el texto reproducido es el del olivo (cap.\ \ref{cap:hacienda}). \\
Ley 1664 (orig.\ 386) & 1936 & Partida «Defensa en los Ríos para Poblaciones», imputada en las obras de 1938 (cap.\ \ref{cap:defensas}). \\
Decretos 17.606 y 17.640 & 03/1934 & Categorías del Registro Civil y de las Receptorías de Rentas (cap.\ \ref{cap:poblacion}). \\
Ley 1269 (orig.\ 11078)* & 27/09/1929 & \textbf{Texto obtenido} ---el Digesto lo reproduce de la \textit{Recopilación General de Leyes} de Gavino Ojeda, y el Boletín de 1929 la publica con ese mismo número original, que sigue la numeración correlativa de leyes y decretos de esos años---. Promulgada el 01/10/1929 (Cornejo--Uriburu); B.O.\ Nº 1.292 del 11/10/1929. Autoriza a la Municipalidad de Campo Santo a reconcentrar el agua del río Mojotoro desde las juntas del Wierna con el La Caldera, y a vigilar y prohibir nuevas bocas-toma no autorizadas por ley. \textbf{No vigente}: el Digesto de Diputados la lista entre las leyes de objeto cumplido, plazo vencido o no vigentes (cap.\ \ref{cap:redes}). \\
Decreto 1966 H* & 16/10/1946 & En acuerdo de ministros: limita a 382 l/s la derivación de la Zona I (ríos Wierna, La Caldera y afluentes) mientras el Wierna lleve menos de 2.500 l/s, y encarga el reparto a la Dirección General de Hidráulica, con colaboración de las municipalidades de Campo Santo y La Caldera. B.O.\ Nº 2.686 (cap.\ \ref{cap:redes}). \\
Acordada 4529/75 & 1975 & Radio de ujiería con límite norte en el río Vaqueros (cap.\ \ref{cap:opacidad}). \\
Decreto 1591/80 & 10/11/1980 & Permiso precario al Club de Regatas Campo Alegre sobre el perilago (cap.\ \ref{cap:tierrafiscal}). \\
Ley 7070 & --- & Protección del ambiente; art.\ 49, audiencia pública (caps.\ \ref{cap:vaqueros} y \ref{cap:loteo}). \\
Ley 7322 y Ley 6994 & sanc.\ 21/10/2004\newline prom.\ Dto.\ 2593/04 & Región Metropolitana de Salta para el transporte de pasajeros, con La Caldera y Vaqueros entre sus municipios; crea la Autoridad Metropolitana del Transporte y SAETA (cap.\ \ref{cap:redes}). \\
Ley 6064 y Ley 7281 & 1983 y --- & Promoción de la actividad turística (caps.\ \ref{cap:poblacion} y \ref{cap:tierrafiscal}). \\
Decreto 2448/85* & 1985 & Reglamenta la Ley 6064: nómina de socios y declaración de bienes en cada solicitud; declaraciones juradas anuales; informes semestrales de la Dirección Provincial de Turismo; pérdida del beneficio por transferencia o arriendo; adhesión municipal eximiendo tasas. B.O.\ Nº 12.356. \\
Ley 6838 & --- & Contrataciones de la provincia; arts.\ 10, 12, 13 y 97 (caps.\ \ref{cap:pdua}, \ref{cap:defensas} y \ref{cap:hacienda}). \\
Ley 5348/78 & --- & Procedimientos administrativos; art.\ 76, rectificación de errores materiales (cap.\ \ref{cap:tierrafiscal}). \\
Ley 3383 & --- & Art.\ 22 inc.\ a: zona de camino de sesenta metros (cap.\ \ref{cap:cot}). \\
Ley nacional 21.477 & --- & Prescripción adquisitiva administrativa a favor del Estado; modificada por la 24.320 (cap.\ \ref{cap:tierrafiscal}). \\
Ley nacional 26.160 & --- & Emergencia territorial indígena y relevamiento (cap.\ \ref{cap:resistencias}). \\
Decreto 397/14 & 11/02/2014 & Coordinador del Plan Urbano Ambiental de Vaqueros; préstamo BID 1855/OC-AR (cap.\ \ref{cap:pdua}). \\
Res.\ 409 D/15 & 08/06/2015 & Contrato con AGEMET S.R.L.\ para el plan de Vaqueros (cap.\ \ref{cap:pdua}). \\
Decreto 3069/14 & 14/10/2014 & Promoción turística, La Caldera S.A.; B.O.\ Nº 19.409. Contrato no publicado: la ficha web no ofrece «Ver anexo» (cap.\ \ref{cap:poblacion}). \\
Decreto 1433/15 & 29/04/2015 & Aprueba el convenio del 10/03/2014 por el que la Provincia cede al municipio su parte del arancel del Ce.N.A.T.\ (36\,\% del formulario, en conjunto), con destino a seguridad vial. Convenio publicado como imagen vinculada; B.O.\ Nº 19.531 (cap.\ \ref{cap:hacienda}). \\
Decreto 463/15 & 09/02/2015 & Promoción turística, Proyecto 1 S.A.; B.O.\ Nº 19.485. Contrato no publicado: «anexo para su consulta en oficinas de nuestra repartición» (cap.\ \ref{cap:poblacion}). \\
Decreto 1968/15 & 08/06/2015 & Prescripción administrativa del terreno de la Escuela Nº 4118; plano 000962 (cap.\ \ref{cap:tierrafiscal}). \\
Decreto 2949/15 & 24/08/2015 & Comodato de 41 ha fiscales a Proyecto 1 S.A.\ (cap.\ \ref{cap:tierrafiscal}). \\
Decreto 3760/15 & 13/11/2015 & Interés provincial de la Fiesta Patronal de La Caldera (cap.\ \ref{cap:opacidad}). \\
Res.\ 59 D/14 & 30/01/2014 & Plano del Club de Campo Las Vertientes (cap.\ \ref{cap:loteo}). \\
Res.\ SRH 251/16, 424--426/16, 205 y 210--212/17, 327/19, 75/2020, 182/2021, 51, 167, 206, 209 y 258/2022, 61/2024 & 2016--2024 & Serie del arroyo Chaile o Teodolinda en Vaqueros, parcela por parcela; la 75/2020 incluye zona inundable (cap.\ \ref{cap:ribera}). \\
Res.\ SRH 94/12, 155/13 y 201/13 & 2012--2013 & Río Castellanos (matrícula 4027, con advertencia de no construir viviendas cerca de la ribera y comunicación a Inmuebles), arroyo Los Nogales (matrícula 1905) y río Vaqueros (sección B) (cap.\ \ref{cap:ribera}). \\
Res.\ SRH 338/16, 23/2021, 125/2021, 107/2022, 57/2025 y 130/2025 & 2016--2025 & Serie del río y el arroyo Vaqueros; la 130/2025 rectifica la 009/12 sobre la matrícula 216 (cap.\ \ref{cap:ribera}). \\
Res.\ SRH 428/15, 304/17, 181/2020, 44/2021, 174/2021, 94/2022 y 23/2024 & 2015--2024 & Serie del río Wierna y del Las Nieves--Wierna; la 44/2021 incluye zona inundable sobre la matrícula 4097 (cap.\ \ref{cap:ribera}). \\
Res.\ SRH 208/13, 292/13 y 142/2025 & 2013 y 2025 & Comisión del arroyo Quintín o Pacará (catastro 3989) y determinación del arroyo Quintín (catastro 3756) que cierra la comisión 292/13 (cap.\ \ref{cap:ribera}). \\
Res.\ SRH 396/11, 418/11, 011/12 y 020/12 & 2011--2012 & Comisión y determinación de la línea de ribera del arroyo Gallinato en Chacras del Gallinato, matrícula 3709; coordenadas no publicadas (cap.\ \ref{cap:loteo}). \\
Res.\ 579 D/15 y 585 D/15 & 10 y 12/08/2015 & Barrio La Misión y Chacras del Gallinato (cap.\ \ref{cap:loteo}). \\
Res.\ 646 D/21 & 26/11/2021 & Plano 1.164 del loteo El Durazno (cap.\ \ref{cap:loteo}). \\
Res.\ Delegada 369 D/22 & 24/05/2022 & Aceptación de la donación sobre la matrícula 6.558 (cap.\ \ref{cap:loteo}). \\
Res.\ SRH 383/14 y serie 2013--2015 & 18/12/2014 & Línea de ribera de los arroyos El Durazno y Guaranguay, y los veinticinco actos de la serie (cap.\ \ref{cap:ribera}). \\
Res.\ SRH 023/14 & 06/02/2014 & A pedido del intendente, designa a Torres y determina la línea de ribera del río La Caldera a lo largo del pueblo; 134 coordenadas en el anexo vinculado (B.O.\ Nº 19.250 y 19.251). Art.\ 3: sujeta la determinación a defensas con gaviones, expte.\ 34-253.934/13 (cap.\ \ref{cap:ribera}). \\
Res.\ SRH 360/16 & 07/11/2016 & Línea de ribera de ambas márgenes del río Wierna, matrículas 4097, 4953, 391, 389, 1268, 3718 y 3438. Expte.\ 34-171.172/15. B.O.\ Nº 19.951. \textbf{Coordenadas no publicadas}, «a disposición en la sede» (cap.\ \ref{cap:ribera}). \\
Res.\ SRH 56/2025, 108/2025, 111/2025 y 16/2026 & 2025--2026 & Determinaciones y rectificación con puntos remitidos a la sede (cap.\ \ref{cap:ribera}). \\ Res.\ SMyE 20/2022 & 2022 & Acuerdo entre Minería y Aguas del Norte sobre la zona de exclusión de canteras junto a las tomas (591 ha). \textbf{No publicada} (caps.\ \ref{cap:amparo} y \ref{cap:aguabaja}). \\ Res.\ SOP 896/23, 314/24, 951/24, 300/26 y 479/26 & 2023--2026 & Serie de las defensas sobre el río La Caldera: legajo (896/23), rescisión (314/24, 19/03/2024), nuevo legajo (951/24, 01/10/2024), legajo de encauzamiento (300/26, 04/06/2026) y rescisión (479/26, 03/09/2026; B.O.\ Nº 22.266 del 10/09/2026) (cap.\ \ref{cap:defensas}). \\
\midrule
\multicolumn{3}{@{}p{13cm}@{}}{\small\itshape Otras normas incorporadas en la etapa final del relevamiento. Las leyes que el Digesto de la Cámara de Diputados lista entre las no vigentes o de objeto cumplido se identifican como tales en cada fila.} \\
Ley 319 (orig.\ 159)* & 28/11/1884 & \textbf{Texto obtenido.} Promulgada el 01/12/1884 (Solá). Autoriza a la Municipalidad de La Caldera a percibir derechos de inhumación (un peso por adulto, cincuenta centavos por párvulo; gratis para los pobres de solemnidad) y la obliga a llevar un libro de defunciones. No vigente. \\
Res.\ MPyDS 36/22* & 15/02/2022 & Manual de misiones y funciones de la Secretaría de Minería y Energía (B.O.\ Nº 21.178). El Programa Gestión y Policía Minera promueve la actividad, controla su cumplimiento y evalúa los informes de impacto; el Subprograma Áridos asesora a los municipios sobre la extracción y la inspecciona (cap.\ \ref{cap:amparo}). \\
Ley 8405* & dic.\ 2023 & Modifica el art.\ 34 del Código de Procedimientos Mineros (Ley 7141). Promulgada por el Decreto 119; publicada el 22/12/2023. Los informes de impacto ambiental mineros se presentan ante la Secretaría de Minería, que los remite al organismo ambiental; sus observaciones no son vinculantes, aunque su rechazo debe fundarse (cap.\ \ref{cap:amparo}). \\
Res.\ SMyE 86/21, 196/22, 134/24 y 176/25* & 2021--2026 & Cuota fija anual de las canteras de tercera categoría: \$810/ha (2022), \$1.945/ha (2023) y \$45.520/ha (desde 2025; Res.\ 134/24 del 23/12/2024, B.O.\ Nº 21.863). La Res.\ 176/25 de la Secretaría de Minería (29/12/2025; B.O.\ Nº 22.121 del 30/01/2026; firma Carrizo) la fija en \$65.320/ha desde 2026, con destino al Fondo Especial de Promoción Minera. \\
Res.\ MEySP 733D/26* & 02/09/2026 & B.O.\ Nº 22.263. Adjudica con fondos del BIRF (Subvención TF0C8847-AR) los proyectos ejecutivos del Nodo Logístico y Parque Industrial de Güemes, con desagües pluviales, por USD 98.271. YSATÍ Ingeniería, en APCA con Aetos, integró la lista corta (cap.\ \ref{cap:amparo}). \\
Ley 8518* & 20/11/2025 & Presupuesto 2026. Decreto 842/25; B.O.\ Nº 22.092 del 15/12/2025. Gastos totales \$3,89 billones. En el departamento: azud del dique Campo Alegre (\$541,5 millones y \$191,1 millones de emergencia) y Ruta 9 Salta--La Caldera (\$18.026,5 millones, DNV). Las defensas del río Yacones y el camino por la margen derecha del río La Caldera figuran sólo entre las obras «con financiamiento a confirmar», sin monto (cap.\ \ref{cap:amparo}). \\
Ley 8511* & 20/11/2025 & \textbf{Ley de ministerios.} Decreto 789/25; B.O.\ Nº 22.080 del 26/11/2025. Deroga las Leyes 8171 y 8274; suprime el Ministerio de Infraestructura; da a la Jefatura de Gabinete las tierras fiscales, el catastro, las obras públicas, la vivienda y las relaciones con los municipios (art.\ 19); reúne en Producción y Minería los recursos hídricos, la minería, la autoridad ambiental y el ordenamiento territorial (art.\ 25); tope de treinta secretarías de Estado (art.\ 30) (caps.\ \ref{cap:ribera} y \ref{cap:tierrafiscal}). \\
Ley 7913* & 01/12/2015 & \textbf{Digesto Jurídico} de la provincia; B.O.\ Nº 19.874 del 03/10/2016. Clasifica las leyes publicadas hasta el 31/07/2015 en seis anexos: vigentes generales permanentes (I), generales transitorias (II) y particulares (III); de objeto cumplido o plazo vencido (IV y V); y derogadas (VI) (apéndice \ref{cap:propuestas}). \\
Ley 5061* & 29/10/1976 & \textbf{Texto obtenido.} B.O.\ Nº 10.105 del 08/11/1976 (Gadea). Deroga las leyes de creación de cuatro municipios: San José de Orquera (Metán, Ley 3292/58), Nuestra Señora de Talavera y Gaona (Anta, Leyes 4030/65 y 4124/66) y San Agustín (Cerrillos, Ley 4829/74), que se reintegran a su jurisdicción anterior. \textbf{No alcanza a la Ley 4365 de Vaqueros.} \\
Ley 5114* & 22/03/1977 & \textbf{Texto obtenido.} Ley de facto (Gadea), B.O.\ Nº 10.213. Crea la Secretaría de Estado de Municipalidades en el Ministerio de Gobierno, Justicia y Educación, en reemplazo de la Dirección General de Coordinación de Municipalidades. La Secretaría aparece, sin mención de esta ley, en el capítulo \ref{cap:politica}. No vigente. \\
Ley 5686* & 20/11/1980 & \textbf{Texto obtenido.} Ley de facto, B.O.\ Nº 11.115. Intendentes y presidentes de comisión designados y removidos por el Ejecutivo, que ejercen también las funciones de los concejos con autorización previa del gobierno; Consejos Asesores Comunales de vecinos, con un mínimo de tres miembros en los municipios de tercera categoría (cap.\ \ref{cap:poblacion}). No vigente. \\
Ley 6334* & 03/09/1985 & \textbf{Texto obtenido.} B.O.\ Nº 12.311. Expropia 2 ha 1.736,80 m² de la Finca Vaqueros (lote II~a, catastro 1.229; dominio de Manuel Serrey y otros) para viviendas; donación al IPDUV con reversión a dos años (cap.\ \ref{cap:expropiacion}). No vigente. \\
Ley 6464* & 14/10/1986 & \textbf{Texto obtenido.} Decreto de necesidad y urgencia 2876/86, tenido por ley por el Decreto 1500/87 al vencer el plazo legislativo; B.O.\ Nº 12.757 del 03/08/1987. Expropia 3.710,08 m² del catastro 1.280 de la Finca Vaqueros, de la sucesión testamentaria de Carlos Serrey, para la planta potabilizadora de Vaqueros (cap.\ \ref{cap:expropiacion}). No vigente. \\
Ley 6895* & 03/09/1996 & \textbf{Texto obtenido.} B.O.\ Nº 15.006. Expropia terrenos del \textbf{departamento Capital} sobre la margen derecha del río Vaqueros (matrículas 91.676 y 75.465) para venderlos a sus ocupantes, y encarga a la Administración General de Aguas las defensas aguas abajo del puente de la Ruta 9. El Decreto 1962/96 vetó la expropiación de la matrícula 91.676 y la parcelación (arts.\ 1 y 4). No vigente. \\
Ley 7460* & 23/08/2007 & \textbf{Texto obtenido.} B.O.\ Nº 17.706. Dona a la Municipalidad de La Caldera el catastro 802, parcela única, manzana 50 del barrio Jardín, para uso público y recreativo. No vigente (objeto cumplido). \\
Ley 7699* & 10/11/2011 & \textbf{Texto obtenido.} B.O.\ Nº 18.728. Autoriza a donar a la Municipalidad de La Caldera la matrícula 1.846, «con el cargo de ser destinado exclusivamente al desarrollo integral de la zona», exenta de gastos. No vigente (objeto cumplido). \\
\bottomrule
\end{longtable}

\section*{Advertencia de uso sobre la Ley 1349}

\begin{cautela}
\textbf{La Ley 1349 está derogada.} El artículo 93 de la Ley 8126 de Régimen de Municipalidades (sancionada el 13/11/2018, B.O.\ 17/12/2018) la derogó junto con las Leyes 3293, 5156, 5814 y 6133. Todo lo que este libro dice de la 1349 vale para el período 1933--2018; desde entonces rige la 8126, que no clasifica a los municipios en categorías, no prevé comisiones municipales y da a todos un Concejo Deliberante y un intendente electo. \textit{Los capítulos \ref{cap:poblacion}, \ref{cap:politica} y \ref{cap:plan} aplican al presente la Ley 8126.}

El texto disponible de la Ley 1349 conserva artículos superados por la reforma constitucional de 1986 y por la legislación electoral posterior. Su artículo 22 dice que los intendentes ``serán nombrados y removidos por el P.E.\ con acuerdo del Senado''; su artículo 6 fija mandatos de concejal de dos años; sus artículos 43 a 50 figuran expresamente derogados por la Ley 6133, y el 31 por la Ley 5814.

Toda cita de esta ley debe contrastarse contra el texto constitucional vigente y advertir la superposición. Hasta su derogación la ley fue un fósil administrativo en uso parcial, y ése es en sí mismo un dato sobre el objeto de este libro. Con una precisión: el art. 22 \textbf{no siempre fue letra muerta}. Rigió hasta la reforma constitucional de 1986, y por eso los primeros intendentes del retorno democrático fueron designados por el gobernador y no electos, según se desarrolla en el capítulo \ref{cap:politica}.
\end{cautela}

\section*{Fuentes estadísticas utilizadas}

Censo Nacional de Población, Hogares y Viviendas 2010 (cuestionario ampliado) y 2022, procesados con Redatam 7 (CEPAL/CELADE), área 66077 y fracciones 66077-01, 02 y 03. Para la desagregación municipal, cuadro \textit{Provincia de Salta. Total de viviendas y de población, según gobierno local. Año 2022} (INDEC), con la codificación del Anexo 1 de la Resolución 144/2022. Cifras definitivas del departamento (cuadro 1.17, por gobierno local): población total 12.299, en viviendas particulares 12.246, en viviendas colectivas 53; viviendas particulares 4.991. La cifra de 12.458 corresponde a los resultados provisorios de 2023 y no se utiliza. Censos Nacionales Agropecuarios 1988 (Anuario Estadístico 2004, cuadro 6.1.1, nivel provincial), 2002 (nivel provincial en el Anuario y departamental en los cuadros de Salta), 2008 (con las cautelas del cap.\ \ref{cap:tierra}) y 2018 (sistema de consultas, cuadros C2 1 y C2 2, nivel departamental).

\section*{Los anexos que faltaban}

\begin{cautela}
El Anexo de la Resolución 82/16 ---Reglamento de Asignación de Lotes de Interés Social--- se obtuvo tarde, y está en el repositorio de anexos del Boletín como imagen escaneada; su artículo 5º, al que remite el plazo de exhibición de padrones, regula la publicidad de los sorteos con un único plazo de diez días de anticipación; por remisión, ése es el mínimo de exhibición (cap.\ \ref{cap:tierrafiscal}). El Anexo de la Ordenanza 437 de Vaqueros ---Reglamento de Fraccionamiento de Tierras---, que contiene las obras exigibles a un loteador en el municipio más presionado del departamento, no está en el Boletín, pero se obtuvo completo del digesto municipal (cap.\ \ref{cap:vaqueros}).
\end{cautela}
